\begin{note}
	The consequence of the above: if \( m = n \), then \( T \) is an isomorphism iff \( M_{\cB_F,\cB_G}(T) \)
	is invertible. Which is also \textcolor{red}{equivalent} to \( \det(M) \) is invertible, being a
	unit: which deducess from:
	\begin{itemize}
		\item \( \det(AB) = \det(A) \cdot \det(B) \).
		\item If \( A^* \) is the classical adjoint of \( A  \implies A A^* = A^* A  = \det(A) \cdot I_n\).
	\end{itemize}
	This is the only difference with a field.
\end{note}
Now what happen if we try to change basis? Actually see that the matrices is dependence on the
choice of basis, just as vector spaces did. Suppose we have another basis of \( F \):
\[
	\tilde{\cB}_F: \tilde{e}_1, \ldots, \tilde{e}_n
\]
one can then write
\[
	\begin{aligned}
		\tilde{e}_j     = \sum_{i=1}^n b_{ij} e_i \implies (b_{ij}) & = M_{\tilde{\cB}_F,\cB_F}(\Id)              \\
		                                                            & =: \textcolor{red}{M_{\tilde{\cB}_F,\cB_F}}
	\end{aligned}
\]
which is the \textcolor{red}{change of basis matrix}, this is of course invertible since \( \Id \)
is invertible. Now suppose \( T: F \to G \) and we shall have the following commutative diagram:
\[
	\begin{tikzcd}
		F(\tilde{\cB}_F) \arrow[r, "T"] \arrow[d, "\Id"'] & G(\tilde{\cB}_G) \arrow[d, "\Id"] \\
		F(\cB_F) \arrow[r, "T"]                 & G(\cB_G)
	\end{tikzcd}
\]
Now utilize \textbf{Equation} \ref{eq:matrix-morphism}, we have
\[
	M_{\tilde{\cB}_F,\tilde{\cB}_G}(T) =  M_{\cB_F,\cB_G}(T) \cdot M_{\tilde{\cB}_G, \cB_G}
	= M_{\tilde{\cB}_G, \cB_G} \cdot M_{\tilde{\cB}_F, \tilde{\cB}_G}(T)
\]
\begin{conclusion}
	\leavevmode
	\begin{equation}
		\label{eq:matrix-morphism-1}
		M_{\tilde{\cB}_F,\tilde{\cB}_G}(T) = (M_{\tilde{\cB}_G, \cB_G})^{-1} \cdot M_{\cB_F,\cB_G}(T) \cdot
		M_{\tilde{\cB}_F, \cB_F}
	\end{equation}
\end{conclusion}

\begin{exercise}
	If \( e_1,\ldots, e_n \) be basis of \( F \) and \( \tilde{e_j} = \sum b_{ij} e_i \), s.t. \(
	(b_{ij}) \) is invertible, we then have \( \tilde{e_1}, \ldots, \tilde{e_n} \) is a basis of \( F \).
\end{exercise}
Suppose now \( F = G\) and \( T: F \to F \) be \( R \)-linear. If we choose the same basis on
\textcolor{blue}{source \& target}, we get the matrix \( M = (a_{ij}) \), and thus \textbf{Equation}
\ref{eq:matrix-morphism-1} tells us if we change the basis, then \( M \) gets
replaced by a \underline{similar matrix} \( M' = B^{-1}M B \), with \( B \) invertible matrix.

\begin{conclusion}
	If \( f \) is free of rank \( n \), then we have:
	\[
		\begin{aligned}
			\qo{\{ R\text{-linear maps } T: V \to V \}}{\text{similarity}}          & \overset{\text{bij}}{\longleftrightarrow}
			\qo{M_n(R)}{\text{similarity}}                                                                                                           \\
			T                                                                       & \rsa M_{\cB, \cB} \text{ w.r.t. some basis } \cB \text{ of } T \\
			R\text{-linear transformation corresp. to } A \text{ w.r.t. some basis} & \leftsquigarrow A
		\end{aligned}
	\]
\end{conclusion}
\underline{Back to our setting}: Let \( \KK \) be a field, and \( V \) be vector space over \( \KK \), with
\( n = \dim_{\KK} V < \infty \). Let \( T: V \to V \) be linear transformation. We saw previously,
we can view \( V \) as a module \( M = (V,T) \) over \( R = \KK[x] \) in \textbf{Proposition}
\ref{prop:module-linear-map}.

\begin{proposition}
	Let \( u_1, \ldots, u_n \) be a \( \KK \)-bases of \( V \) and let \( A = (a_{ij}) \) the matrix
	of \( T \) w.r.t. this basis. Let \( M = (V,T) \). Then we have \textcolor{red}{a short exact sequence} of \( \KK[x] \)-module as
	follows:
	\[
		0 \to \KK[x]^n \xrightarrow{\vp} \KK[x]^n \xrightarrow{\psi} M \to 0
	\]
	where
	\begin{enumerate}
		\item \( \psi(e_i) = u_i \) for \( 1\leq i \leq n \) where \( e_1,\ldots, e_n \) be standard
		      basis of \( \KK[x]^n \).
		\item The matrix of \( \vp \) w.r.t. the standard basis is \( x\cdot I_n - A \) and \( x \) is
		      the formal variable in \( \KK[x] \).
	\end{enumerate}
	\begin{note}
		Recall that \( T(u) := x\cdot u \).
	\end{note}
\end{proposition}

\begin{proof}
	Since \( u_1,\ldots, u_n \) generates \( M \) over \( \KK \), thus they also generates it over \(
	\KK[x] \implies \psi\) is surjective. Now we have
	\[
		\begin{aligned}
			\psi \circ \vp (e_j)                                         & = \psi\bace{xe_j - \sum_{i=1}^n a_{ij}e_j } = T(u_j) - \sum_{i=1}^n
			a_{ij} u_j = 0 \text{\todo{by def of matrix}}                                                                                                                                             \\
			\implies \; \exists \;    \overline{\psi}: \text{Coker}(\vp) & \to M \text{ s.t. }
			\overline{\psi}(\overline{u}) = \psi(u) \; \forall \; u \in \KK[x]^n                                                                                                                      \\
			                                                             & \text{where } \psi(xe_j)                                             = \underbrace{x\underbrace{\psi(e_j)}_{u_j}}_{T(u_j)}
		\end{aligned}
	\]
	thus \( \im(\vp) \subseteq \ker(\psi)  \).
	To show that
	\[
		\ker(\psi) = \im (\vp)
	\]
	We already know that \( \psi \) is surjective with \(
	\dim_{\KK} M = n \). To prove it is an isomorphism of \( \KK \)-vector space, hence of \( R
	\)-modules, it is enough to show:
	\[
		\dim_{\KK} \frac{\KK[x]^n}{\im(\vp)} \leq n
	\]
	\begin{note}
		We already seen that by first isomorphism theorem:
		\[
			\frac{\KK[x]^n}{\ker(\vp)} \cong \im(\psi) = M
		\]
		and thus if
		\[
			\dim_{\KK} \frac{\KK[x]^n}{\im(\vp))} \leq  \dim_{\KK} M = n
		\]
		\[
			\implies \ker(\psi) \subseteq \im(\vp) \implies \ker(\psi) = \im(\vp)
		\]
	\end{note}
	Hence it is enough to show: \( \overline{e_1}, \ldots, \overline{e_n} \) generate this over \( \KK \).
	It is clear that \( \{x^j \overline{e_i} \; | \; 1 \leq i \leq m, j \geq 0\} \) generates \(
	\frac{\KK[x]^n}{\im(\vp)} \) over \( \KK \). Thus it is enough to show by induction on \( j \geq 0 \) that
	\[
		x^j \overline{e_i} \in \langle \overline{e_1}, \ldots, \overline{e_i} \rangle =  \sum_{l=1}^n \KK
		\overline{e_l}
	\]
	The \textcolor{red}{key point} is that in \( \frac{\KK[x]^n}{\im(\vp)} \) we have
	\[
		\begin{aligned}
			\bace{x\overline{e_i} - \sum_{l=1}^n a_{li}\overline{e_l}} & \in \im(\vp)                                                                \\
			\implies x\overline{e_i}                                   & \in \sum_{l=1}^n \KK \overline{e_l}                                         \\
			\implies x^j \overline{e_i}                                & \in \sum_{l=1}^n \KK x^{j-1}\overline{e_l} \subseteq \sum_{l=1}^n
			\KK \overline{e_l}                                                                                                                       \\
			                                                           & \text{where the last inclusion given by \textcolor{blue}{induction on $j$}}
		\end{aligned}
	\]
	This completes the proof that \( \im(\vp) = \ker(\psi) \). Let's now show that \( \vp \) is
	\textcolor{blue}{injective}. We have
	\[
		R^n \xrightarrow{\vp} R^n \xrightarrow{\psi} M \to 0
	\]
	with \( \ker(\psi) = \im(\vp) \implies \)
	\[
		\begin{aligned}
			M                                   & \cong \qo{R^n}{\im(\vp)} \\
			\implies \underbrace{\rank(M)}_{=0} & = n - \rank(\im(\vp))    \\
			\implies \rank(\im(\vp))            & = n                      \\
		\end{aligned}
	\]
	thus
	\[
		\begin{aligned}
			0                          \to \ker(\vp) \to R^n \xrightarrow{\vp} & \im(\vp) \to 0 \\
			\qo{R^n}{\ker(\vp)}                                                & \cong \im(\vp) \\
			\implies \rank(\ker(\vp))                                          & = 0            \\
		\end{aligned}
	\]
	since \( \ker(\vp) \subseteq  \) free \( R \)-module \( \implies \ker(\vp) = 0 \).
\end{proof}
We have that
\[
	\vp: \KK[x]^n \hookrightarrow \KK[x]^n
\]
the general theorem on submodules of free modules over PIDs tells us that there exists \( A, B_1,
B_2 \in
M_n(\KK[x]) \) being invertible, s.t.
\[
	B_1 \cdot (x \cdot I_n - A ) \cdot B_2 =
	\begin{pmatrix}
		1 &        &   &     &   &        & 0 \\
		  & \ddots &   &     &   &        &   \\
		  &        & 1 &     &   &        &   \\
		  &        &   & f_i &   &        &   \\
		  &        &   &     & 1 &        &   \\
		  &        &   &     &   & \ddots &   \\
		0 &        &   &     &   &        & 1 \\
	\end{pmatrix}
\]
s.t. \( f_1,\ldots, f_r \) are monic polynomials with degree \( geq 1 \) and \( f_1 \mid f_2 \mid
\cdots \mid f_r \).
